%==============================================================================
\chapter{Analysis and correction theory}\label{ch:anatheory}
%==============================================================================

This appendix gives the signal processing behind the Analysis view of Chapter~\ref{ch:analysis}, in the order the data flows. Each \gterm{tf}{transfer function} is estimated, the positions are aligned and averaged, the average is smoothed, the correction is designed and optionally aligned to a subwoofer, and the result is rendered as an \gterm{fir}{FIR}.

\section{Welch estimation}\label{sec:thwelch}
For each file the plugin estimates the loudspeaker-room \gterm{tf}{transfer function} $H(f)=Y(f)/X(f)$, where $x$ is the emitted stimulus (channel~1 of the file) and $y$ the captured microphone signal (channel~2). If a microphone calibration is in effect, it is divided out of $H$ at this stage (eq.~\eqref{eq:miccal}).

\subsection{The $H_1$ estimator}
A direct bin-by-bin division $Y/X$ is unusable in a room. It is dominated by measurement noise wherever $X$ is small, and a single FFT gives an estimate with $100\,\%$ standard error. SuperMoTo instead uses the classical \gterm{welch}{Welch} averaged-periodogram approach~\cite{Welch1967}. The recording is cut into overlapping segments, each is windowed and transformed, and the auto- and cross-spectra are averaged over the $L$ segments,
\begin{equation}
  \hat P_{xx}(f) = \frac{1}{L}\sum_{\ell=1}^{L} |X_\ell(f)|^2,
  \qquad
  \hat P_{xy}(f) = \frac{1}{L}\sum_{\ell=1}^{L} X_\ell^{*}(f)\,Y_\ell(f).
\end{equation}
The transfer function is then formed as the \gterm{h1}{$H_1$ estimator}
\begin{equation}
  \hat H(f) = \frac{\hat P_{xy}(f)}{\hat P_{xx}(f)} .
  \label{eq:h1}
\end{equation}
$H_1$ is the least-squares estimator that is optimal when the noise is on the output (the microphone path), the realistic case for an acoustic measurement, and it rejects that noise as more segments are averaged. (Its companion $H_2=P_{yy}/P_{yx}$ assumes input noise. Both bracket the true response, and they coincide only when the \gterm{coherence}{coherence} is unity.) The derivation and noise properties of $H_1/H_2$ are standard material in \cite{BendatPiersol,ShinHammond,OppenheimSchafer}.

\subsection{Windowing, overlap and resolution}
Each segment is weighted by a Hann window before the FFT, to control spectral leakage~\cite{Harris1978}. Consecutive segments overlap by 50\,\%, the standard choice that makes the Hann-windowed segments use the recording almost uniformly (Welch's WOSA scheme). The window length $W$, from 16384 to 262144 samples, sets the frequency resolution
\begin{equation}
  \Delta f = \frac{f_s}{W},
\end{equation}
about \SI{0.67}{\hertz} at $f_s=\SI{44.1}{\kHz}$ with the default $W=65536$. Bins where $\hat P_{xx}$ is essentially zero are \gterm{regularization}{regularised} by a small term $\varepsilon = 10^{-10}\max_k \hat P_{xx} + 10^{-30}$ added to the denominator of~\eqref{eq:h1}, so silent bins do not blow up.

\section{Delay removal and complex averaging}\label{sec:thavg}
Each microphone position is at a different distance from the speaker, so each $\hat H$ carries a different bulk propagation delay, a linear phase term $e^{-j2\pi f \tau}$. Averaging the complex responses without removing it would let the position-dependent phase ramps cancel each other and destroy the result. The plugin therefore estimates the delay $\tau$ of each measurement from its \gterm{ir}{impulse response} (inverse FFT of $\hat H$), taken as the position of the peak of $|h(t)|$. A peak in the second half of the buffer is interpreted as a small negative delay that wrapped around. The linear phase is then removed,
\begin{equation}
  \hat H'(f) = \hat H(f)\, e^{+j 2\pi f \tau},
\end{equation}
which time-shifts every measurement so its main arrival sits at $t=0$. This is the excess-delay normalisation used in swept-sine room measurement practice~\cite{MullerMassarani}. Appendix~\ref{ch:delays} discusses what this arrival-time estimate measures, and where it fails.

The delay-aligned responses of all $N$ positions are then averaged, and the magnitude and the phase of that average come from different means, for reasons that are worth separating.

The \textbf{magnitude} is a power mean,
\begin{equation}
  M(f) = \sqrt{\frac{1}{N}\sum_{n=1}^{N} \bigl|\hat H'_n(f)\bigr|^2 },
  \label{eq:powermean}
\end{equation}
because $|\overline{H}|$ is not $\overline{|H|}$. Above the room's transition frequency the positions no longer agree on phase, and the modulus of a complex mean of $N$ near-random phasors reads about $1/\sqrt N$ low. For instance, on a ten-position set, one gets \SI{0.2}{\dB} at \SI{200}{\hertz} but \SI{12}{\dB} at \SI{13}{\kilo\hertz}, at bins where every individual curve is flat. Inverting that would boost a cancellation that exists at no microphone position. Equation~\ref{eq:powermean} is the spatial average of the sound field, and averaging out the position-dependent \gterm{comb}{comb filtering} caused by reflections is what makes several positions worthwhile in the first place~\cite{Toole,Mourjopoulos1994}.

The \textbf{phase} comes from the complex sum $S(f)=\sum_n \hat H'_n(f)$, provided that this sum is statistically meaningful. The validity and reliability of this sum are quantified by the vector coherence:
\begin{equation}
  R(f) = \frac{|S(f)|}{N\,M(f)} \in [0,1],
  \label{eq:veccoherence}
\end{equation}
which is unity where every position agrees and zero where they are independent. $N$ unit phasors that agree on nothing still sum to $\sqrt N$, so $R$ has a floor of $1/\sqrt N$ that means nothing. Removing it gives the weight actually used,
\begin{equation}
  w(f) = \operatorname{clamp}\!\left( \frac{N R^2(f) - 1}{N-1},\, 0,\, 1 \right),
  \label{eq:phaseweight}
\end{equation}
so $w=0$ reads "no better agreement than chance". The average is then the power-mean magnitude carried by a direction blended between the one the positions voted for and the minimum-phase direction of that same magnitude,
\begin{equation}
  \bar H(f) = M(f)\,\frac{b(f)}{|b(f)|},\qquad
  b(f) = w\,\frac{S}{|S|} + (1-w)\,e^{\,j\varphi_{\min}(f)},
  \label{eq:blendedaverage}
\end{equation}
with $\varphi_{\min}$ the minimum-phase argument of $M$, obtained from the real cepstrum as in Section~\ref{sec:minphase}. The two are blended as phasors on the short arc, as the subwoofer alignment blends its own, so no phase is ever unwrapped and nothing can jump.

\paragraph{Why the phase needs a weight at all.} Where the positions disagree, $S$ nearly cancels, and its argument is then the direction of a residual between near-random phasors. The phase can shift by \SI{180}{\degree} between adjacent frequency bins even as the magnitude $M$ varies smoothly. Because this abrupt phase flip leaves the magnitude unaffected, it cannot be detected by analyzing the magnitude spectrum alone. However, a \gterm{linphase}{linear-phase} rendering exposes this transition by forcing the complex trajectory to pass through a near-zero point on the unit circle. Measured on a ten-position campaign whose right loudspeaker had $R=0.07$ near \SI{670}{\hertz}, the exported filter carried a \SI{-25}{\dB} notch of $Q\approx500$ at \SI{697}{\hertz}, audible as a ringing tone, where the minimum-phase rendering of the same design (using the magnitude alone) was \SI{+2.3}{\dB} at that frequency and showed nothing. Applying~\eqref{eq:blendedaverage} leaves the magnitude correction completely unchanged at \SI{0.000}{\dB} RMS over the \SIrange{200}{10000}{\hertz} range. Meanwhile, the crossover delay estimate shifts by merely \SI{0.02}{\milli\second} and the summation efficiency by \SIrange{0.02}{0.04}{\dB}, while the notch is entirely eliminated: no bin between \SI{650}{\hertz} and \SI{750}{\hertz} falls below \SI{-0.4}{\dB}, as in the minimum-phase rendering. This approach is also fundamentally justified: spatial variations in phase do not reflect the acoustic properties of the overall listening area. Attempting to correct these position-dependent phase differences might optimize the response for a single seat, but it will inevitably degrade performance for the rest of the audience.

\paragraph{What a value of $R$ is worth.} $R$ is the mean resultant length of circular statistics, and for a wrapped-normal or von~Mises spread of phases it fixes the scatter directly: the circular standard deviation is $\sigma=\sqrt{-2\ln R}$. That is what Table~\ref{tab:coherence} reads out, beside the weight~\eqref{eq:phaseweight} gives at $N=10$. The two agree where they should: at $R=1/\sqrt{10}=0.316$, the value $N$ independent positions produce on their own, the scatter is \SI{87}{\degree} (phases spread almost uniformly around the circle) and the weight has fallen to $0.003$. The bias correction in~\eqref{eq:phaseweight} is the standard small-sample correction for the squared mean resultant length~\cite{Kutil2012}.

\begin{table}[ht]
\centering
\renewcommand{\arraystretch}{1.2}
\begin{tabular}{lrrrrrrr}
\toprule
$R$                        & 0.99 & 0.95 & 0.90 & 0.80 & 0.70 & 0.50 & 0.32\\
circular sd $\sigma$ (deg) &    8 &   18 &   26 &   38 &   48 &   68 &   87\\
weight $w$, $N=10$         & 0.98 & 0.89 & 0.79 & 0.60 & 0.43 & 0.17 & 0.00\\
\bottomrule
\end{tabular}
\caption{What the positions' agreement is worth. $R$ is the vector coherence
of~\eqref{eq:veccoherence}, $\sigma=\sqrt{-2\ln R}$ the phase scatter it
implies, and $w$ the weight~\eqref{eq:phaseweight} gives that phase in the
average. Independent positions put a floor of $1/\sqrt N$ under $R$ --- 0.58 at
three positions, 0.32 at ten, 0.22 at twenty --- which is what the correction
removes.}
\label{tab:coherence}
\end{table}

\paragraph{What it looks like on a real set.} Table~\ref{tab:measuredcoherence} is $R$ measured per octave on a ten-position campaign, both main loudspeakers. It follows the shape Appendix~\ref{ch:spatial} predicts from the correlation length: near unity through the bass, where a \SI{30}{\centi\metre} cluster samples a single correlation cell, and falling through the midrange as the positions begin to sample independently. The weight therefore lets the average's phase through in full below a few hundred hertz and fades it out above, without anyone choosing a frequency. That settles which phase is information. Whether inverting it is worth what it costs is a separate question, and Section~\ref{sec:thphaselimit} answers it with a frequency after all.

\begin{table}[ht]
\centering
\renewcommand{\arraystretch}{1.2}
\small
\begin{tabular}{lrrrrrrrrr}
\toprule
octave from (Hz) & 31 & 63 & 125 & 250 & 500 & 1k & 2k & 4k & 8k\\
\midrule
left main   & 0.98 & 0.99 & 0.90 & 0.89 & 0.79 & 0.67 & 0.50 & 0.67 & 0.46\\
right main  & 0.97 & 0.99 & 0.90 & 0.88 & 0.78 & 0.68 & 0.59 & 0.59 & 0.34\\
\bottomrule
\end{tabular}
\caption{Vector coherence per octave band, labelled by the band's lower edge,
on ten positions spanning \SI{45}{\centi\metre} by \SI{34}{\centi\metre}
with four of them \SI{6}{\centi\metre} higher. The floor for ten independent
positions is 0.32.}
\label{tab:measuredcoherence}
\end{table}

\paragraph{How rare the failure is, and why that is the problem.} On that same campaign the defect occupied \emph{one bin}. Taking the largest phase step between neighbouring bins anywhere in the analysis band, the left loudspeaker reached \SI{11.9}{\degree} and the right \SI{154.9}{\degree}, the latter at \SI{697.8}{\hertz} and at exactly one of the \num{29632} bins in band. With~\eqref{eq:blendedaverage} in place they fall to \SI{8.3}{\degree} and \SI{9.7}{\degree}, both at the crossover where the subwoofer alignment turns the phase on purpose, and no bin anywhere exceeds \SI{90}{\degree}. A defect that occupies one bin in thirty thousand, leaves the magnitude untouched and is plainly audible is exactly the kind that a curve on a screen will not show and a listening test will.

\section{Synchronized swept-sine deconvolution}\label{sec:thsweep}
With the sweep stimulus, the \gterm{tf}{transfer function} can be estimated by deconvolution instead of by \gterm{welch}{Welch} averaging. The implementation follows the synchronized swept-sine formulation of Novak, Lotton and Simon~\cite{Novak2015}.

\subsection{The synchronized sweep}\label{sec:syncsweep}
The stimulus is the exponential (logarithmic) sweep
\begin{equation}
  x(t) = \sin\!\left( 2\pi f_1 L \left( e^{t/L} - 1 \right) \right),
  \label{eq:sweep}
\end{equation}
from $f_1$ to $f_2$, where $f_2$ is the Nyquist frequency and $f_1$ lies a whole number of octaves below it, near \SI{11}{\hertz} at every sample rate (Appendix~\ref{ch:sweepband} explains why). Its instantaneous frequency is $f(t) = f_1 e^{t/L}$, so the sweep rate $L$ is the time the stimulus takes to multiply its frequency by $e$, and the sweep reaches frequency $f$ at
\begin{equation}
  t(f) = L \ln \frac{f}{f_1}.
  \label{eq:sweeptime}
\end{equation}

What makes the sweep synchronized is that $L$ is not taken directly from the requested duration $T$, but quantized so that $f_1 L$ is an integer:
\begin{equation}
  L = \frac{1}{f_1}\,
      \operatorname{round}\!\left( \frac{f_1 T}{\ln (f_2/f_1)} \right).
  \label{eq:sweepL}
\end{equation}
The actual sweep duration is then $L \ln (f_2/f_1)$, slightly different from the $T$ requested.

\subsection{Deconvolution by the analytic inverse}
The recorded response is deconvolved by the sweep's analytic spectral inverse, the closed form obtained from the stationary-phase approximation of the sweep's spectrum~\cite{Novak2015},
\begin{equation}
  X^{-1}(f) = 2\sqrt{\frac{f}{L}}\;
    \exp\!\left( -\,j\,2\pi f L
      \left( 1 - \ln\frac{f}{f_1} \right) + j\,\frac{\pi}{4} \right).
  \label{eq:sweepinv}
\end{equation}
The recording is transformed, multiplied by~\eqref{eq:sweepinv}, and transformed back. Because the inverse is known in closed form, nothing is divided by a measured spectrum. There is no ill-conditioned denominator to \gterm{regularization}{regularise} (contrast the $\varepsilon$ term in~\eqref{eq:h1}), and the whole band, magnitude and phase, comes from a single sweep rather than from an average over segments.

\subsection{Why the harmonic orders separate}\label{sec:sweepharm}
A distortion product of order $m$ at frequency $f$ is generated while the sweep is passing $f/m$. By~\eqref{eq:sweeptime} that happens
\begin{equation}
  \Delta t_m = L \ln m
  \label{eq:sweepharmoffset}
\end{equation}
seconds before the linear response at the same frequency. After deconvolution, therefore, the \gterm{ir}{impulse responses} of the different orders sit at different places on the time axis. The linear impulse response lands at the system's propagation delay, and the order-$m$ impulse response lands $L \ln m$ seconds before it, wrapping circularly towards the end of the buffer.

This is where the quantization~\eqref{eq:sweepL} earns its place. With $f_1 L$ an integer, the sweep's phase at every arrival time $\Delta t_m$ is an exact multiple of $2\pi$, so the harmonic impulse responses emerge not merely separated in time but with their true phase. Deconvolving a recording whose $L$ was not quantized (a file from another tool, say) still separates the orders in time, but their phase relationship is lost. Under the Welch estimator, by contrast, the distortion products are simply part of the measured output and end up folded into the linear estimate.

\subsection{Extracting the linear and harmonic responses}
The deconvolution uses an FFT the next power of two at least as long as the recording, and never shorter than the sweep, so the wrapped harmonic responses land far from the linear one. The first $W$ samples of the deconvolved response hold the linear impulse response. They are transformed onto exactly the same $W/2+1$ bin grid the Welch path produces, so everything downstream is method-agnostic.

Harmonic orders \numrange{2}{5} are extracted, each centred $L \ln m$ seconds before the linear peak. The extraction half-width is bounded by $0.45$ times the spacing to the next order, which shrinks as $L \ln\!\big(m/(m-1)\big)$, so neighbouring orders cannot leak into one another. An order is dropped when the sweep is too short to separate it cleanly. Across microphone positions the harmonic magnitudes are averaged in power, because distortion phase is not coherent from one position to the next and the vector average of Section~\ref{sec:thavg} would cancel the distortion instead of measuring it.

\section{Fractional-octave smoothing}\label{sec:thsmooth}
The $1/n$-octave \gterm{smoothing}{smoothing} is a complex moving average, magnitude and phase together, over a logarithmic band $[\,k/r,\;k\,r\,]$ around each bin $k$, with $r = 2^{\text{fraction}/2}$, evaluated in $O(n)$ with a prefix sum. Working on a constant-$Q$ grid follows Hatziantoniou and Mourjopoulos~\cite{Hatziantoniou2000}. The $1/2$-octave setting fills the large perceptual gap between $1/3$ and $1$~octave.

With frequency-dependent smoothing, the octave fraction used at bin $k$ is interpolated log-linearly in frequency between the low- and high-frequency settings,
\begin{equation}
  \mathrm{frac}(f) = \mathrm{frac}_{\mathrm{LF}}
    + t(f)\,\bigl(\mathrm{frac}_{\mathrm{HF}}-\mathrm{frac}_{\mathrm{LF}}\bigr),
  \quad
  t(f) = \mathrm{clamp}\!\left(\frac{\log_2 f - \log_2 f_{\mathrm{lo}}}
                                     {\log_2 f_{\mathrm{hi}} - \log_2 f_{\mathrm{lo}}},\,0,\,1\right),
\end{equation}
with anchors $f_{\mathrm{lo}}=\SI{100}{\hertz}$ and $f_{\mathrm{hi}}=\SI{10}{\kHz}$. This is distinct from the constant-$Q$ widening that any fractional-octave window already has. That widening grows the window in \si{\hertz} as frequency rises but keeps the octave fraction constant, whereas here the octave fraction itself varies with frequency.

\section{Correction design}\label{sec:thcorrection}
Starting from the normalised smoothed average $h = \bar H_s / g_{\mathrm{ref}}$, with $g_{\mathrm{ref}}$ the \SIrange{200}{2000}{\hertz} mean magnitude, the ideal inverse is the complex reciprocal
\begin{equation}
  C_0(f) = \frac{1}{h(f)} = \frac{\overline{h(f)}}{|h(f)|^{2}} ,
  \label{eq:inverse}
\end{equation}
which flattens magnitude and zeroes phase at once. Four mechanisms \gterm{regularization}{regularise} and shape $C_0$ into the correction that is actually applied.

\paragraph{1. Soft-knee boost ceiling.} Rather than hard-clamping the boost, which leaves a kink at every notch, the boost portion of the correction magnitude in dB is bent towards a ceiling $L$ (default $+12\,\si{\dB}$) with a $\tanh$ knee of width $\kappa=\min(6,L)\,\si{\dB}$:
\begin{equation}
  C_{\mathrm{dB}} \;\longmapsto\;
  \begin{cases}
    C_{\mathrm{dB}}, & C_{\mathrm{dB}} \le L-\kappa,\\[2pt]
    (L-\kappa) + \kappa\,\tanh\!\dfrac{C_{\mathrm{dB}}-(L-\kappa)}{\kappa}, & C_{\mathrm{dB}} > L-\kappa.
  \end{cases}
\end{equation}
The boost passes untouched up to $L-\kappa$, then saturates smoothly and asymptotes to $L$. Cuts ($|C|<1$) are left alone. This is the practical form of regularised inversion discussed in~\cite{Kirkeby1999}.

\paragraph{2. Low-frequency slope.} $C$ is multiplied by a 2nd-order Butterworth high-pass target at $f_c=\SI{30}{\hertz}$,
\begin{equation}
  H_{\mathrm{hp}}(f) = \frac{(j\omega)^2}{(j\omega)^2 + \sqrt{2}\,(j\omega) + 1},
  \qquad \omega = f/f_c ,
\end{equation}
so below $f_c$ the correction rolls off rather than chasing the woofer's natural roll-off to infinity.

\paragraph{3. Correction level.} The correction level $\lambda$ morphs continuously from bypass to full correction by raising the correction to a power in the log/phase domain,
\begin{equation}
  C \;\longmapsto\; |C|^{\lambda}\, e^{\,j\lambda\arg C},
  \qquad \lambda \in [0,1],
\end{equation}
so $\lambda=0$ gives unity and $\lambda=1$ the full inverse.

\paragraph{4. Analysis band.} Outside the user range $[f_{\mathrm{lo}},f_{\mathrm{hi}}]$ the correction is faded back to unity through a half-octave raised-cosine skirt $w(f)$,
\begin{equation}
  C \;\longmapsto\; \bigl(1-w(f)\bigr) + w(f)\,C .
\end{equation}

\section{Subwoofer phase alignment}\label{sec:thsubalign}
\paragraph{Preserving the relative timing.} The subwoofer set is taken at the same positions as the main and paired by load order. Each subwoofer measurement is delay-anchored on the corresponding main measurement's delay rather than having its own delay estimated and removed. If each file were independently zeroed, the physically meaningful main-vs-sub arrival-time difference at each position would be discarded. Anchoring keeps it, so the averaged subwoofer response $\bar S(f)$ retains its true phase relative to $\bar H$.

\paragraph{The all-pass alignment term.} From the unit-magnitude phasor of the subwoofer, $\hat s$, an \gterm{allpass}{all-pass} factor is built and blended towards "no change" above the \gterm{crossover}{crossover},
\begin{equation}
  A(f) = \frac{(1-w_x) + w_x\,\hat s(f)}{\bigl|(1-w_x) + w_x\,\hat s(f)\bigr|},
  \qquad C \longmapsto C\cdot A(f),
\end{equation}
where the alignment weight $w_x(f)$ is 1 at and below the crossover and is released to 0 over one octave above it (raised cosine). The renormalisation makes $A$ strictly phase-only, leaving the main's magnitude correction unchanged. Blending the phasors, rather than scaling an unwrapped phase angle, keeps the rotation on the short arc between $1$ and $\hat s$, so $A$ applies only the minimal phase shift and never unwinds the subwoofer's large accumulated low-frequency phase. The \gterm{polarity}{polarity} inversion negates $\bar S$.

\paragraph{Bulk time alignment.} The subwoofer phasor entering the blend is advanced by the Mains delay $t_a$ before its phase is taken,
\begin{equation}
  \hat s(f) = \frac{\hat S(f)}{|\hat S(f)|},\qquad
  \hat S(f) = \pm\,\bar S_s(f)\,e^{\,j 2\pi f\, t_a},
\end{equation}
so when $t_a$ matches the measured offset the corrected main's phase is already flat through the crossover and $A\to 1$. The recommended $t_a$ comes from a power-weighted least-squares fit of the unwrapped, delay-anchored subwoofer phase over $[\,f_c/2,\,2f_c\,]$, the local \gterm{groupdelay}{group delay} $\hat\tau=-\,\mathrm{d}\varphi/\mathrm{d}\omega$. Appendix~\ref{ch:delays} develops this estimator and the condition the applied delays must satisfy.

\section{The complete correction filter}\label{sec:fullfilter}
Collecting the stages of Sections~\ref{sec:thcorrection} and~\ref{sec:thsubalign}, the correction applied at each frequency bin is the single complex expression
\begin{equation}
  \boxed{\;
  C(f) = \bigl(1-w(f)\bigr)
       + w(f)\,A(f)\,
         \mathcal{P}_{\lambda}\!\Bigl[\,
           H_{\mathrm{hp}}(f)\;\mathcal{B}\!\bigl[\,1/h(f)\,\bigr]
         \Bigr]\; }
  \label{eq:fullfilter}
\end{equation}
evaluated on the \gterm{welch}{Welch} bin grid ($f_k = k\,f_s/W$). Section~\ref{sec:threnders} resamples and inverse-transforms this spectrum $C(f)$, not an analytic filter, into the exported \gterm{fir}{FIR}. Reading~\eqref{eq:fullfilter} from the inside out, the operators are exactly the design steps above:
\begin{align}
  h(f) &= \frac{\bar H_s(f)}{g_{\mathrm{ref}}}
        && \text{normalised smoothed average,}\\[2pt]
  \tfrac{1}{h(f)} &= \frac{\overline{h(f)}}{|h(f)|^{2}}
        && \text{exact complex inverse, eq.~\eqref{eq:inverse},}\\[2pt]
  \mathcal{B}[z] &= z\,10^{\,\bigl(\ell(z_{\mathrm{dB}})-z_{\mathrm{dB}}\bigr)/20},\quad
        z_{\mathrm{dB}}=20\log_{10}|z|
        && \text{soft-knee boost limit (phase kept),}\\[2pt]
  \ell(g) &= \begin{cases}
               g, & g \le L-\kappa,\\
               (L-\kappa) + \kappa\tanh\!\dfrac{g-(L-\kappa)}{\kappa}, & g > L-\kappa,
             \end{cases}
        && \kappa=\min(6,L),\ L=\text{max boost (dB)},\\[2pt]
  H_{\mathrm{hp}}(f) &= \frac{(j\omega)^2}{(j\omega)^2+\sqrt{2}\,(j\omega)+1},\quad
        \omega=\frac{f}{f_c}
        && \text{LF slope, } f_c=\SI{30}{\hertz},\\[2pt]
  \mathcal{P}_{\lambda}[z] &= |z|^{\lambda}\,e^{\,j\lambda\arg z}
        && \text{correction level } \lambda\in[0,1],\\[2pt]
  \hat s(f) &= \frac{\pm\,\bar S_s(f)\,e^{\,j2\pi f t_a}}{\bigl|\bar S_s(f)\bigr|}
        && \text{time-aligned sub phasor, } t_a=\text{Mains delay,}\\[2pt]
  A(f) &= \frac{(1-w_x)+w_x\,\hat s(f)}{\bigl|(1-w_x)+w_x\,\hat s(f)\bigr|}
        && \text{sub phase alignment } (A\equiv 1 \text{ if no sub}),\\[2pt]
  w(f) &\in [0,1]
        && \text{analysis-band weight (raised-cosine skirts).}
\end{align}
Two boundary cases complete the definition. Where $|h(f)|$ is negligible, the inverse is replaced by the boost ceiling rather than diverging. Cuts ($z_{\mathrm{dB}}\le 0$) pass $\mathcal{B}$ unchanged since $L-\kappa\ge 0$, so only boosts are limited. Setting $\lambda=1$, no subwoofer ($A\equiv1$) and a fully open analysis band ($w\equiv1$) reduces~\eqref{eq:fullfilter} to the boost-limited, LF-sloped exact inverse $H_{\mathrm{hp}}\,\mathcal{B}[1/h]$. The linear-phase export then applies one more operator to $C$, the phase limit of Section~\ref{sec:thphaselimit}.

\paragraph{The ceiling limits a magnitude, and only a magnitude.} $\mathcal{B}$ bounds $|1/h|$. It does nothing to $\arg(1/h)=-\arg h$, which turns fastest exactly where $|h|$ is smallest. A capped magnitude carried by an uncapped phase is no longer the inverse of anything, and the mismatch grows with how hard the ceiling is working. Two things keep it in proportion. The phase weight of~\eqref{eq:phaseweight} has already dealt with the case that matters most, a null the positions disagree about, where the phase was never information to begin with, and a null they all agree about is a property of the listening area whose phase is worth inverting as far as the ceiling lets the magnitude follow it, within the crossover region where the phase is kept at all (Section~\ref{sec:thphaselimit}). In the specific configuration analyzed above, the boost ceiling is never reached. The unconstrained inverse filter requires a maximum boost of only \SI{5.1}{\dB} and \SI{6.4}{\dB} for the two main channels, well below the \SI{12}{\dB} limit. Consequently, the bounding operator $\mathcal{B}$ acts as an identity function across all in-band frequency bins. Conversely, when a system does encounter its boost ceiling, the calculated correction phase becomes increasingly unreliable inside deep notches where the filter declines to apply further boost. The appropriate remedy for such deep nulls is not to increase the filter length (which merely implements an unrealistic target with higher precision) but rather to apply broader frequency smoothing or to lower the boost ceiling itself.

\section{Limiting the phase to the crossover region}\label{sec:thphaselimit}
Equation~\eqref{eq:fullfilter} inverts the average's phase wherever the weight~\eqref{eq:phaseweight} lets it through. With "Limit phase" on, the default, the linear-phase export keeps that designed phase only where the subwoofer alignment needs it and hands it over to the minimum-phase phase of the same correction above,
\begin{equation}
  \Lambda[C](f) = |C(f)|\,
     \frac{w_p(f)\,\hat c(f) + \bigl(1-w_p(f)\bigr)\,\hat m(f)}
          {\bigl|\,w_p(f)\,\hat c(f) + \bigl(1-w_p(f)\bigr)\,\hat m(f)\,\bigr|},
  \qquad \hat c = \frac{C}{|C|},
  \label{eq:phaselimit}
\end{equation}
where $\hat m(f)$ is the unit phasor of the minimum-phase response with magnitude $|C|$, built by the cepstrum of Section~\ref{sec:minphase}, and
\begin{equation}
  w_p(f) = \begin{cases}
             1, & f \le 2f_x,\\[2pt]
             \tfrac12 + \tfrac12\cos\!\bigl(\pi\log_2\tfrac{f}{2f_x}\bigr), & 2f_x < f < 4f_x,\\[2pt]
             0, & f \ge 4f_x,
           \end{cases}
  \label{eq:phaselimitweight}
\end{equation}
with $f_x$ the crossover. The magnitude does not move, and from $4f_x$ up $\Lambda[C]$ is exactly the filter the minimum-phase rendering exports. The blend runs on the short arc, like those of~\eqref{eq:blendedaverage} and the alignment term.

\paragraph{Why stop at twice the crossover.} The alignment weight $w_x$ has fallen to zero at $2f_x$, so above it the all-pass $A$ is the identity and the only phase left to invert is the room's. Below it, the room's phase has to stay. $A$ steers the main onto the subwoofer on the assumption that the main's own phase is already flat there. Replacing the average's phase with minimum phase at every frequency cost \SI{1}{\dB} of summation efficiency in the offline model of a ten-position campaign ($-1.6$ against $-0.6$). Limiting it at $2f_x$ cost nothing: below $2f_x$ the limited and unlimited exports are the same filter, and measured through the crossover band they differ by \SI{0.01}{\dB}.

\paragraph{Why above it at all.} Only the part of the average's phase that every seat shares can be inverted everywhere~\cite{NeelyAllen, Mourjopoulos1994}. The rest stays at each seat, and a linear-phase filter puts half of whatever it does not cancel \emph{ahead} of the direct sound, where the ear masks little of it: a \gterm{preecho}{pre-echo}. At $R\approx0.8$--$0.9$, the coherence of the \SIrange{250}{1000}{\hertz} octaves in Table~\ref{tab:measuredcoherence}, the weight still passes most of the phase while the positions scatter by \SIrange{26}{38}{\degree}. That residual is exactly what was heard, as a short pre-reverberation before impacts, on a campaign re-exported after the fix of Section~\ref{sec:thavg}. Making the weight stricter does not cure it: raising $w$ to the fourth power still leaves the \SI{500}{\hertz} band \SI{10}{\dB} above minimum phase. Positions that agree this well can still leave each seat enough to hear.

Table~\ref{tab:phaselimit} measures it through the whole system: from \SIrange{500}{8000}{\hertz} the unlimited export put \SIrange{14}{24}{\dB} more energy than the limited one in the \SI{30}{\milli\second} before the direct sound, and \SIrange{14}{30}{\dB} more in the last \SI{3}{\milli\second}. The limited one is within \SI{2}{\dB} of minimum phase over the longer window. At \SI{250}{\hertz}, inside the fade, it removes \SIrange{3}{10}{\dB}.

Nor did the inversion buy anything there. The arrival of each octave's envelope was measured against the direct sound at every seat, through the whole system. From \SIrange{500}{4000}{\hertz} the median seat was time-coherent to within \SI{0.2}{\milli\second} in every octave with no correction at all and through every export alike, and the worst seat to within \SI{0.3}{\milli\second} through any export. Against minimum phase, the unlimited export moved none of the octaves from \SI{250}{\hertz} up by more than \SI{0.11}{\milli\second}. Only below $2f_x$, where the alignment works, did linear phase move an arrival: the \SI{125}{\hertz} octave \SIrange{0.4}{0.5}{\milli\second} earlier and the \SI{63}{\hertz} one \SI{0.9}{\milli\second} later, the main being steered onto the subwoofer.

\begin{table}[ht]
\centering
\small
\setlength{\tabcolsep}{4.5pt}
\renewcommand{\arraystretch}{1.2}
\begin{tabular}{lrrrrrr}
\toprule
 & \SI{250}{\hertz} & \SI{500}{\hertz} & \SI{1}{\kilo\hertz} & \SI{2}{\kilo\hertz} & \SI{4}{\kilo\hertz} & \SI{8}{\kilo\hertz} \\
\midrule
\multicolumn{7}{l}{\textit{\SIrange{30}{3}{\milli\second} before the direct sound}}\\
minimum phase       & $-49/{-50}$ & $-57/{-57}$ & $-64/{-66}$ & $-68/{-67}$ & $-68/{-68}$ & $-67/{-67}$\\
linear, limited     & $-35/{-34}$ & $-56/{-56}$ & $-64/{-64}$ & $-67/{-67}$ & $-68/{-68}$ & $-67/{-67}$\\
linear, not limited & $-32/{-24}$ & $-38/{-34}$ & $-45/{-42}$ & $-46/{-47}$ & $-44/{-48}$ & $-50/{-53}$\\
noise floor         & $-68$ & $-71$ & $-72$ & $-71$ & $-68$ & $-67$\\
\midrule
\multicolumn{7}{l}{\textit{\SIrange{3}{1}{\milli\second} before}}\\
minimum phase       & $-62/{-65}$ & $-56/{-59}$ & $-64/{-63}$ & $-59/{-56}$ & $-74/{-69}$ & $-77/{-61}$\\
linear, limited     & $-42/{-44}$ & $-56/{-58}$ & $-64/{-63}$ & $-59/{-56}$ & $-69/{-69}$ & $-73/{-61}$\\
linear, not limited & $-39/{-26}$ & $-28/{-28}$ & $-36/{-33}$ & $-40/{-42}$ & $-39/{-39}$ & $-43/{-47}$\\
\bottomrule
\end{tabular}
\caption{What limiting the phase does, measured through the whole system at ten
positions on a campaign crossed at \SI{70}{\hertz}, 8192 taps, left/right main.
Energy arriving ahead of the direct sound, per octave, relative to that octave's
energy in the \SI{300}{\milli\second} after it, in dB, through causal band
filters so the analysis cannot move anything earlier; power mean over the
positions. The noise floor is the same \SI{27}{\milli\second} window taken
\SIrange{400}{150}{\milli\second} ahead of the direct sound, beyond the reach of
any of the filters. Every response is deconvolved from its sweep whole: Welch
averaging of the same sweeps reads some \SI{20}{\dB} more ahead of the direct
sound in the \SIrange{250}{1000}{\hertz} octaves, in every row alike.}
\label{tab:phaselimit}
\end{table}

\paragraph{Why the limit is applied to the finished correction.} The minimum-phase phasor $\hat m$ in~\eqref{eq:phaselimit} is that of $|C|$, not of the average. The average's minimum phase also carries what the correction never inverts: the band edges of the measurement, and the loudspeaker's roll-off below the analysis range, where $w$ returns $C$ to unity. Inverting their phase without their magnitude is a pre-echo of its own. On a synthetic set of three positions agreeing on a \SI{1}{\kilo\hertz} all-pass, a limit built from the average left \SI{-5.7}{\dB} of the render's energy ahead of its peak, more than no limit at all ($-13$). Built from $|C|$, as~\eqref{eq:phaselimit} does, it leaves \SI{-38}{\dB} and matches the minimum-phase design to \SI{0.001}{\degree} above $4f_x$. The test suite checks both: the match above $4f_x$, and a pre-echo at least \SI{20}{\dB} below the unlimited design's.

\paragraph{What it leaves.} Below $2f_x$ the linear-phase filter still puts energy ahead of the direct sound. On the campaign of Table~\ref{tab:phaselimit} the \SI{63}{\hertz} octave carries \SI{-14}{\dB} and the \SI{125}{\hertz} one \SIrange{-16}{-17}{\dB} within \SI{30}{\milli\second}, against \SI{-20}{\dB} and \SI{-26}{\dB} in minimum phase. That is the alignment itself, and the design delay sets its size: crossed at \SI{100}{\hertz}, the \SI{125}{\hertz} octave carried \SI{-4}{\dB} with the estimated delay and \SIrange{-13}{-16}{\dB} with one \SI{14}{\milli\second} shorter, the summation being the same at both to \SI{0.03}{\dB}. What it buys is a summation that does not depend on that delay: measured, \SIrange{0.2}{0.5}{\dB} more level through the crossover band than minimum phase at the estimated delay and \SI{2.7}{\dB} more at the shorter one; predicted, between \SI{-0.24}{\dB} and \SI{-0.97}{\dB} of summation efficiency for any delay from \SIrange{10}{40}{\milli\second}, a range over which minimum phase falls as low as \SI{-8.4}{\dB}. Whether that low-frequency energy is audible is a listening question. \SI{30}{\milli\second} is less than two periods at \SI{63}{\hertz}.

\section{Rendering the FIR}\label{sec:threnders}
Both exports turn a complex spectrum into a finite \gterm{ir}{impulse response} on the \gterm{fir}{FIR}'s frequency grid, with $N=\text{nextPow2}(\text{firLength})$ bins.

\subsection{Linear- and mixed-phase rendering}
The resampled spectrum is given a half-length linear phase shift, multiplying bin $k$ by $e^{-j\pi k}=(-1)^k$, before the inverse FFT. Because the correction is a mixed-phase filter, deliberately carrying phase and \gterm{allpass}{all-pass} content including the subwoofer alignment, it is not causal on its own, and centring its energy at sample $N/2$ is what makes it realisable as an ordinary FIR. A Tukey window (10\,\% taper) centred on $N/2$ then cleans up the edges. FIR design by frequency sampling and windowing is standard~\cite{OppenheimSchafer}.

The cost of centring is a constant latency of $\approx N/2$ samples on that output. SuperMoTo's per-output convolution uses the WDL zero-latency (uniformly partitioned) convolution engine~\cite{Gardner1995}, so the filter adds no algorithmic block latency beyond that fixed $N/2$ offset.

\subsection{Minimum-phase rendering}\label{sec:minphase}
As defined in Appendix~\ref{ch:minphasedef}, a \gterm{minphase}{minimum-phase} response is the unique causal, stable response whose log-magnitude and phase form a \gterm{hilbert}{Hilbert-transform} pair (eq.~\eqref{eq:hilbert}), so it is fully determined by $|C(f)|$~\cite{OppenheimSchafer}. SuperMoTo computes it by the real-\gterm{cepstrum}{cepstrum} method. With $c[n]=\mathcal{F}^{-1}\{\log|C|\}$ the real cepstrum, the minimum-phase log-spectrum is recovered by keeping its causal part,
\begin{equation}
  \hat c[n] = c[n]\,\ell[n], \qquad
  \ell[n] = \begin{cases}
              1, & n=0,\ n=N/2,\\
              2, & 1\le n < N/2,\\
              0, & N/2 < n \le N-1,
            \end{cases}
\end{equation}
and the impulse response is $\mathcal{F}^{-1}\bigl\{\exp \mathcal{F}\{\hat c\}\bigr\}$, tapered only at the tail since its energy is front-loaded near $n=0$~\cite{OppenheimSchafer}. Only $|C|$ enters this construction, and the all-pass term $A$ has unit magnitude, so the subwoofer alignment has no effect on a minimum-phase export.
